\section{The Tengwar}
Phi can be written in either of two ways: using the standard Latin alphabet or the Tengwar script. Tengwar was invented by J. R. R. Tolkien and first described in Appendix E of the third book in his trilogy, \emph{The Lord of the Rings}. There are numerous sources of information available concerning the Tengwar so details and history of the script will not be explored in this document.

The look and feel of the Tengwar characters (\emph{tengwa}) and its diacritical marks (\emph{tehta}) match the ideals and philosophies behind this language. Phi has a small set of consonants and vowels so it is relatively straightforward to map its graphemes to Tengwar characters. The Unicode font \emph{Tengwar Telcontar}\footnote{\hspace{0.1em}\url{http://freetengwar.sourceforge.net/tengtelc.html}} is used in this manual to render Phi text in Tengwar script. A custom keyboard layout\footnote{\hspace{0.1em}\url{http://freetengwar.sourceforge.net/keylayouts.html}} was employed to facilitate typing Tengwar characters.

Tables \ref{tab:consonant-graphemes} and \ref{tab:vowel-graphemes} outline the mapping between the graphemes used in Phi and the corresponding Tengwar characters, as well as the keypresses necessary to produce the characters using the font and the keyboard layout mentioned previously. An application or utility that supports Graphite fonts\footnote{\hspace{0.1em}\url{https://scripts.sil.org/cms/scripts/page.php?site\_id=projects\&item\_id=graphite\_fonts}} is required for the font to render correctly. A link to details of this requirement can be found in the first footnote below.

Written Phi uses the Tengwar’s Quenya mode. Vowel diacriticals are placed on the consonant that precedes the vowel. Since all Phi words start with a consonant, there is never a need to use an initial vowel carrier. Words in Phi often contain two consecutive vowels, but a vowel carrier is still unnecessary: the second vowel is placed below the same consonant. The letters would then be read as consonant, upper vowel, then lower vowel. This and the use of a single Tengwar character for consonant doubles such as ph or sh cause Phi script to be very compact and efficient.
\clearpage
The Phi word for “love” is \emph{lothea}. In the Tengwar script it looks like this:

\bigskip
\leftskip=0.66in
\ttl  \normalfont
\bigskip

\leftskip=0in
This six-letter word in latin characters requires only two \emph{tengwa} to represent it. Many words in Phi are two to three letters long. This occurs mainly in parts of speech such as particles, prepositions, or pronouns. These words are almost always a single \emph{tengwa} with diacriticals in Tengwar script. All of this results in a very concise writing system.

This is the phrase “I love you” in its simplest vernacular form in Phi:

\bigskip
\leftskip=0.66in
\ttl  \normalfont
\bigskip

\leftskip=0in
That is a beautiful sentiment rendered in a beautiful script. In more formal communication in Phi, particles are used to mark parts of speech in order to reduce ambiguity. As well, word markers are often used to distinguish multisyllabic words for clarity. The Phi Tengwar punctuation is outlined in Table 5. As an example, the phrase “I love you” would be written formally as:

\bigskip
\leftskip=0.66in
\ttl ⸱⸱⸱⸱⸱ \normalfont
\bigskip

\clearpage
\begin{table}[H]
  \caption {\label{tab:consonant-graphemes}
    Phi consonant grapheme to Tengwar mapping}
  \renewcommand\arraystretch{2.5}
  \setlength\extrarowheight{-3pt}
  \begin{tabular}{|c|c|c|}
    \hline \rowcolor{greyhead}
    Grapheme &
    Tengwa &
    Keypress \\
    \hline
    p & \tts  & p \\
    \hline \rowcolor{greyrow}
    t & \tts  & t \\
    \hline
    m & \tts  & m \\
    \hline \rowcolor{greyrow}
    n & \tts  & n \\
    \hline
    r & \tts  & shift-r \\
    \hline \rowcolor{greyrow}
    s & \tts  & s \\
    \hline
    s & \tts  & shift-s \\
    \hline \rowcolor{greyrow}
    h & \tts  & h \\
    \hline
    w & \tts  & w \\
    \hline \rowcolor{greyrow}
    l & \tts  & l \\
    \hline
    ph & \tts  & shift-p \\
    \hline \rowcolor{greyrow}
    th & \tts  & shift-t \\
    \hline
    sh & \tts  & shift-c \\
    \hline \rowcolor{greyrow}
    wh & \tts  & shift-k \\
    \hline
  \end{tabular}
\end{table}
\noindent
\begin{table}[H]
  \caption {\label{tab:vowel-graphemes}
    Phi vowel grapheme to Tengwar mapping}
  \renewcommand\arraystretch{2.5}
  \setlength\extrarowheight{-3pt}
  \begin{tabular}{|c|c|c|}
    \hline \rowcolor{greyhead}
    Grapheme & \hspace{1em}
    Tehta \hspace{1em} &
    Keypress \\
    \hline
    i & \tts   & i / option-i\footnotemark \\
    \hline \rowcolor{greyrow}
    e & \tts   & e / option-e \\
    \hline
    o & \tts   & o / option-o \\
    \hline \rowcolor{greyrow}
    a & \tts   & a / option-a \\
    \hline
  \end{tabular}
\end{table}
\footnotetext{\hspace{0.1em}Keypresses for the low diacriticals are given for macOS. Use \emph{ctrl-alt} on Windows instead of \emph{option}.}
\noindent
\begin{table}[H]
  \caption {\label{punctuation}
    Phi punctuation to Tengwar mapping}
  \renewcommand\arraystretch{2.5}
  \setlength\extrarowheight{-3pt}
  \begin{tabular}{|c|c|c|}
    \hline \rowcolor{greyhead}
    Punctuation &
    Tengwar &
    Keypress \\
    \hline
    space & \tts ⸱ & , \\
    \hline \rowcolor{greyrow}
    period & \tts : & shift-, \\
    \hline
  \end{tabular}
\end{table}
